% page 628 du fac-simile (image p-627 du scan)
% bloc 157.5 x 229.0 mm ; marge gauche 8.0 mm
\pgc{-12.700mm}{0.000mm}{
\l{\cel{6}620}
\l{}
\l{vizitar. (\sou{trans.}) Irar renkontre che ulu, pro devo di poli-}
\l{\cel{3}teso, di deferenco. - DEFIRS.}
\l{}
\l{\sou{voco}. I. (a) Sono de la aero ekpulsita de la pulmoni, kande}
\l{\cel{3}ta aero trairas laringo, qua igas ica vibrar.\cel{2}(b) Sono}
\l{\cel{3}artikulata, parolo. (c) Sono modulata segun la skalo mu-}
\l{\cel{3}zikala, aptigita pri kantado. - II. Expreso di la opinio-}
\l{\cel{3}no di persono lor voto. - III. (\sou{metaf.}) (\sou{gram.}) Formo di}
\l{\cel{3}la konjugo, qua indikas se la efiko quan expresas la verbo}
\l{\cel{3}esas agata o subisata da la subjekto.- EFIS.}
\l{}
\l{\sou{vodevilo}. Teatrajo mixita kun kupleti. - DEFIR.}
\l{}
\l{\sou{vokar}. (\sou{trans.}) Invitar (ulu) pri venar, per pronuncar lua}
\l{\cel{3}nomo, o per signalo. EIL.}
\l{}
\l{\sou{vokalo}. Sonuno simpla, emisata per lasar la aero ekirar de}
\l{\cel{3}laringo e trairar la boko sen obstaklo.--II. Singla ek la}
\l{\cel{3}sonuni diversa qui genitesas tale, segun la maniero quale}
\l{\cel{3}onu apertas sua boko por lasar li ekirar.(la vokali : a,}
\l{\cel{3}e, i, o, u). - DEFIS.}
\l{}
\l{\sou{vokalizar}. (\sou{netrans.}) (\sou{muziko)}. I.Exercar su voce, per kantar}
\l{\cel{3}la noto sen pronuncar lua nomo e sen vortizar la sonuno.}
\l{\cel{3}- II. Adjuntar traiti, ornivi, en kantajo. - DEFIRS.}
\l{}
\l{\sou{vokativo}. (\sou{gram.}) Un ek la kazi di la deklino (\sou{che ula lingui}}
\l{\cel{3}\sou{nacionala}) quan onu uzas kande onu su turnas a persono od}
\l{\cel{3}a kozo quan onu egardas quale personon. - DEFIS.}
\l{}
\l{\sou{volar}. (\sou{trans.}) I. Rezolvar agor o ne agor (ulo). - II. Prak-}
\l{\cel{3}tikar plu o min energioze sua povo pri rezolvar koncerne}
\l{\cel{3}agar (ulo) od igar ulu agar (ulo). - III. Rezolvar ferme,}
\l{\cel{3}obtenor (ulo) por su o por altru. - DEFIRS.}
\l{}
\l{\sou{volatila}. (\sou{aludante substanco solida o liquida}). Qua vapores-}
\l{\cel{3}kas facile. - EFIS.}
\l{}
\l{\sou{volfo}. (\sou{zool.}) Quadripedo karnavida, kun muzelo longa e nigra,}
\l{\cel{3}oreli rekta, pilaro flavatra e nigre-buntita, kaudo tufa-}
\l{\cel{3}tra. - L.\cel{1}\sou{canis lupus}.-DE.}
\l{}
\l{\sou{volkano}. I. Abismo subtera, (normale) interne di monto, qua}
\l{\cel{3}lansas, per aperturo (kratero), flami, fumuro, cindri, e}
\l{\cel{3}materii fuzura. - II. (\sou{Metaf.}) (a)\cel{1}\sou{Esar sur volkano}\cel{1}:}
\l{\cel{3}en stando politikala quan karakterizas la minaco da revo-}
\l{\cel{3}luciono. (b)\cel{1}\sou{Viro volkana}\cel{1}: di qua la karaktero esas im-}
\l{\cel{3}petuoza. - DEFIRS.}
\l{}
\l{volontario. Qua armeanesas, qua su engajis pri armeanesko,}
\l{\cel{3}sen obligeso pri lo. - DEFIRS.}
}
